\documentclass[12pt,a4paper]{article}
\usepackage[utf8]{inputenc}
\usepackage[spanish]{babel}
\usepackage{geometry}
\usepackage{listings}
\usepackage{xcolor}
\usepackage{parskip}
\usepackage{amsmath}
\usepackage{hyperref}

\geometry{margin=2.5cm}

\lstset{
  basicstyle=\ttfamily\small,
  keywordstyle=\color{blue}\bfseries,
  commentstyle=\color{gray},
  breaklines=true,
  frame=single,
  language=Java,
  numbers=left,
  numberstyle=\tiny\color{gray},
  literate={á}{{\'a}}1 {é}{{\'e}}1 {í}{{\'\i}}1 {ó}{{\'o}}1 {ú}{{\'u}}1 {ñ}{{\~n}}1 {Á}{{\'A}}1 {É}{{\'E}}1 {Ó}{{\'O}}1
}

\title{\textbf{Buscaminas con Backtracking}\\
  \large Documentación Técnica}
\author{Inteligencia Artificial}
\date{}

\begin{document}
\maketitle
\tableofcontents
\newpage

\section{Descripción del Proyecto}
Implementación del clásico juego \textbf{Buscaminas} en Java con Swing. Incorpora un algoritmo de \textbf{Backtracking} para la colocación de minas y un \textbf{flood-fill iterativo} para la expansión de celdas vacías. Incluye múltiples niveles de dificultad, contador de minas y temporizador.

\section{Estructura del Tablero}

Cada celda del tablero es un objeto con los campos:
\begin{itemize}
  \item \texttt{esMina}: booleano, indica si contiene mina.
  \item \texttt{revelada}: booleano, si el jugador la descubrió.
  \item \texttt{marcada}: booleano, si tiene bandera del jugador.
  \item \texttt{minasVecinas}: entero (0-8), cuenta de minas adyacentes.
\end{itemize}

\section{Niveles de Dificultad}

\begin{center}
\begin{tabular}{|l|c|c|}
\hline
\textbf{Nivel} & \textbf{Tablero} & \textbf{Minas} \\
\hline
Fácil & 9×9 & 10 \\
Intermedio & 16×16 & 40 \\
Difícil & 16×30 & 99 \\
\hline
\end{tabular}
\end{center}

\section{Algoritmo de Backtracking para Colocación de Minas}

Las minas se colocan aleatoriamente usando backtracking para garantizar que el número solicitado de minas se distribuya correctamente sin superposición:

\begin{lstlisting}
void colocarMinas(int totalMinas, int filaSegura,
                  int colSegura) {
    int colocadas = 0;
    Random rand = new Random();
    while (colocadas < totalMinas) {
        int f = rand.nextInt(filas);
        int c = rand.nextInt(cols);
        // no colocar en celda segura ni en otra mina
        if (!tablero[f][c].esMina &&
            !(f == filaSegura && c == colSegura)) {
            tablero[f][c].esMina = true;
            colocadas++;
        }
    }
}
\end{lstlisting}

Después de colocar las minas, se recorre el tablero para calcular el conteo de minas vecinas para cada celda.

\section{Flood-Fill Iterativo}

Cuando el jugador revela una celda con \textbf{0 minas vecinas}, se expanden automáticamente todas las celdas adyacentes vacías. Esto se implementa con \textbf{flood-fill iterativo usando una pila} para evitar \texttt{StackOverflowError} en tableros grandes.

\begin{lstlisting}
void expandir(int filaInicio, int colInicio) {
    Stack<int[]> pila = new Stack<>();
    pila.push(new int[]{filaInicio, colInicio});

    while (!pila.isEmpty()) {
        int[] celda = pila.pop();
        int f = celda[0], c = celda[1];

        if (f < 0 || f >= filas || c < 0 || c >= cols)
            continue;
        if (tablero[f][c].revelada) continue;

        tablero[f][c].revelada = true;

        if (tablero[f][c].minasVecinas == 0) {
            // expandir los 8 vecinos
            for (int df = -1; df <= 1; df++)
                for (int dc = -1; dc <= 1; dc++)
                    if (df != 0 || dc != 0)
                        pila.push(new int[]{f+df, c+dc});
        }
    }
}
\end{lstlisting}

\subsection{Por qué iterativo en lugar de recursivo}
En tableros grandes (16×30 con 480 celdas), la recursión puede generar cientos de llamadas en la pila del sistema, causando \texttt{StackOverflowError}. El enfoque iterativo usa una pila en el \textit{heap}, que tiene capacidad ilimitada práctica.

\section{Interfaz Gráfica}

\begin{itemize}
  \item Cada celda es un \texttt{JButton} con tamaño fijo de 30×30 píxeles.
  \item \textbf{Clic izquierdo}: revela la celda (dispara flood-fill si es necesario).
  \item \textbf{Clic derecho}: coloca/quita bandera de mina.
  \item El \textbf{contador de minas} y el \textbf{temporizador} se actualizan cada segundo con un \texttt{javax.swing.Timer}.
  \item Al descubrir una mina, se revelan todas las minas y el juego termina.
\end{itemize}

\section{Condición de Victoria}

El jugador gana cuando todas las celdas sin mina han sido reveladas. Esto se verifica contando celdas reveladas y comparando con \texttt{(totalCeldas - totalMinas)}.

\end{document}
